% =============================================================================
% Vibe Coding Sandbox - Tool Documentation Template
% =============================================================================
% Usage:
%   1. Copy this file to your tool's docs directory
%   2. Fill in the sections below
%   3. Compile with: pdflatex document.tex
%   4. Or use pandoc: pandoc README.md -o document.pdf
% =============================================================================

\documentclass[11pt,a4paper]{report}

% ── Packages ────────────────────────────────────────────────────────────────
\usepackage[utf8]{inputenc}
\usepackage[T1]{fontenc}
\usepackage{lmodern}
\usepackage[margin=1in]{geometry}
\usepackage{graphicx}
\usepackage{hyperref}
\usepackage{listings}
\usepackage{xcolor}
\usepackage{booktabs}
\usepackage{longtable}
\usepackage{enumitem}
\usepackage{fancyhdr}
\usepackage{titlesec}

% ── Code Highlighting ───────────────────────────────────────────────────────
\definecolor{codebg}{RGB}{248,248,248}
\definecolor{codeborder}{RGB}{221,221,221}
\definecolor{comment}{RGB}{0,128,0}
\definecolor{keyword}{RGB}{0,0,255}
\definecolor{string}{RGB}{163,21,21}

\lstset{
    backgroundcolor=\color{codebg},
    frame=single,
    rulecolor=\color{codeborder},
    numbers=left,
    numberstyle=\tiny\color{gray},
    basicstyle=\ttfamily\small,
    breaklines=true,
    showstringspaces=false,
    commentstyle=\color{comment},
    keywordstyle=\color{keyword}\bfseries,
    stringstyle=\color{string},
    language=bash,
    tabsize=2,
    xleftmargin=2em,
    framexleftmargin=1.5em
}

% ── Hyperlinks ──────────────────────────────────────────────────────────────
\hypersetup{
    colorlinks=true,
    linkcolor=blue,
    filecolor=magenta,
    urlcolor=cyan,
    pdftitle={Tool Documentation},
    pdfauthor={Vibe Coding Sandbox}
}

% ── Header/Footer ───────────────────────────────────────────────────────────
\pagestyle{fancy}
\fancyhf{}
\fancyhead[L]{\leftmark}
\fancyhead[R]{\thepage}
\renewcommand{\headrulewidth}{0.4pt}

% ── Section Formatting ──────────────────────────────────────────────────────
\titleformat{\chapter}[display]
    {\normalfont\huge\bfseries}{\chaptertitlename\ \thechapter}{20pt}{\Huge}
\titleformat{\section}
    {\normalfont\Large\bfseries}{\thesection}{1em}{}
\titleformat{\subsection}
    {\normalfont\large\bfseries}{\thesubsection}{1em}{}

% =============================================================================
% DOCUMENT START
% =============================================================================
\begin{document}

% ── Title Page ──────────────────────────────────────────────────────────────
\begin{titlepage}
    \centering
    \vspace*{3cm}
    {\Huge\bfseries Tool Name\\[0.5cm]}
    {\Large Technical Documentation\\[2cm]}
    {\large\itshape Vibe Coding Sandbox\\[0.3cm]}
    {\large Renoise Themes Project\\[3cm]}
    \vfill
    {\large Generated: \today}
\end{titlepage}

% ── Table of Contents ───────────────────────────────────────────────────────
\tableofcontents
\newpage

% =============================================================================
% CHAPTER 1: OVERVIEW
% =============================================================================
\chapter{Overview}

\section{What This Tool Does}
\textit{[Write a 2-3 sentence summary of the tool's purpose and value proposition.]}

\section{Problem Statement}
\textit{[Describe the problem this tool solves. Why did you build it?]}

\section{Key Features}
\begin{itemize}[leftmargin=2em]
    \item \textit{[Feature 1]}
    \item \textit{[Feature 2]}
    \item \textit{[Feature 3]}
\end{itemize}

% =============================================================================
% CHAPTER 2: ARCHITECTURE
% =============================================================================
\chapter{Architecture}

\section{System Diagram}
\textit{[Include an ASCII diagram or reference an image file.]}

\begin{lstlisting}[language=bash,caption=High-level command]
$ tool-name --input file.txt --output result.json
\end{lstlisting}

\section{Component Breakdown}
\textit{[Describe each module/file and its responsibility.]}

\begin{longtable}{@{}llp{8cm}@{}}
\toprule
\textbf{File} & \textbf{Type} & \textbf{Purpose} \\
\midrule
\endhead
\texttt{lib/parser.js} & Module & Parses input files \\
\texttt{tools/build.sh} & Script & Automation pipeline \\
\texttt{app.js} & Entry & Main application \\
\bottomrule
\end{longtable}

% =============================================================================
% CHAPTER 3: USAGE
% =============================================================================
\chapter{Usage}

\section{Installation}
\begin{lstlisting}[language=bash,caption=Install dependencies]
$ npm install
# or
$ pip install -r requirements.txt
\end{lstlisting}

\section{Running the Tool}
\begin{lstlisting}[language=bash,caption=Basic usage]
$ node tool-name.js [options] <input>
\end{lstlisting}

\section{Configuration}
\textit{[Document environment variables, config files, etc.]}

\begin{longtable}{@{}llp{6cm}@{}}
\toprule
\textbf{Variable} & \textbf{Default} & \textbf{Description} \\
\midrule
\texttt{PORT} & 3000 & Server port \\
\texttt{NODE\_ENV} & development & Runtime environment \\
\bottomrule
\end{longtable}

% =============================================================================
% CHAPTER 4: API / INTERFACE
% =============================================================================
\chapter{API / Interface}

\section{Endpoints / Commands}
\textit{[Document each API endpoint or CLI command.]}

\subsection{Command: \texttt{generate}}
\begin{lstlisting}[language=bash]
$ tool-name generate --format json --output out/
\end{lstlisting}

\textbf{Parameters:}
\begin{itemize}[leftmargin=2em]
    \item \texttt{-f, --format} — Output format (json, xml, csv)
    \item \texttt{-o, --output} — Output directory
\end{itemize}

% =============================================================================
% CHAPTER 5: IMPLEMENTATION DETAILS
% =============================================================================
\chapter{Implementation Details}

\section{Algorithms}
\textit{[Describe any non-trivial algorithms or data structures.]}

\section{Performance Characteristics}
\textit{[Time/space complexity, benchmarks, etc.]}

\section{Known Limitations}
\textit{[Be honest about edge cases, TODOs, and future work.]}

% =============================================================================
% CHAPTER 6: EXAMPLES
% =============================================================================
\chapter{Examples}

\section{Example 1: Basic Workflow}
\textit{[Step-by-step example with expected output.]}

\section{Example 2: Advanced Usage}
\textit{[Show off power-user features.]}

% =============================================================================
% REFERENCES
% =============================================================================
\chapter{References}

\section{File Index}
\begin{longtable}{@{}lp{10cm}@{}}
\toprule
\textbf{File} & \textbf{Description} \\
\midrule
\texttt{README.md} & Main project readme \\
\texttt{CLAUDE.md} & AI assistant context \\
\bottomrule
\end{longtable}

\section{Related Documentation}
\begin{itemize}[leftmargin=2em]
    \item \textit{[Link to related tool docs]}
    \item \textit{[Link to external references]}
\end{itemize}

\end{document}
